\documentclass[11pt,a4paper]{article}
\usepackage[margin=1in]{geometry}
\usepackage{amsmath,amssymb,amsthm}
\usepackage{algorithm}
\usepackage{algpseudocode}
\usepackage{booktabs}
\usepackage{hyperref}

\theoremstyle{definition}
\newtheorem{definition}{Definition}[section]
\newtheorem{theorem}{Theorem}[section]
\newtheorem{lemma}[theorem]{Lemma}
\newtheorem{remark}{Remark}[section]

\title{\textbf{Rigorous Mathematical Specification, Robust Estimation Dynamics, and Subgradient Analysis of Continuous Regression Loss Functions}}
\author{Team Scaibu \\ \small Email: \texttt{jaydeep.9664920749@gmail.com} \\ \small Architecture and Core Engine Team}
\date{October 2026}

\begin{document}
\maketitle

\begin{abstract}
This document presents the complete theoretical formulation and algorithmic specification of continuous regression objective functions: Mean Squared Error (MSE), Mean Absolute Error (MAE), Smooth $L_1$ / Huber Loss, and Asymmetric Quantile (Pinball) Loss. We formalize likelihood estimation under Gaussian and Laplace noise models, derive analytical subgradients, establish robust M-estimator asymptotic properties, provide the complete algorithmic pseudo-code matching the reference implementation, derive computational complexity, calculate a worked numerical example, and address numerical bounds.
\end{abstract}

\section{Notation and Mathematical Preliminaries}

Let $\hat{y} \in \mathbb{R}^N$ and $y \in \mathbb{R}^N$ denote predicted continuous outputs and ground-truth targets respectively across $N \in \mathbb{N}_{\ge 1}$ batch samples.

\begin{itemize}
    \item $e_i \triangleq \hat{y}_i - y_i \in \mathbb{R}$: Residual prediction error for sample $i \in \{1, \dots, N\}$.
    \item $\delta > 0$: Huber loss threshold transitioning between quadratic and linear penalty (default $1.0$).
    \item $\tau \in (0, 1)$: Quantile parameter for asymmetric pinball regression (default $0.5$).
    \item $\mathcal{L}(\hat{y}, y) \in \mathbb{R}$: Scalar loss value after reduction ($\text{reduction} \in \{\text{mean}, \text{sum}\}$).
    \item $g_i \triangleq \frac{\partial \mathcal{L}}{\partial \hat{y}_i} \in \mathbb{R}$: Analytical loss gradient with respect to prediction $\hat{y}_i$.
\end{itemize}

\begin{table}[h]
\centering
\caption{Summary of Mathematical Symbols for Regression Losses}
\begin{tabular}{lll}
\toprule
\textbf{Symbol} & \textbf{Type / Domain} & \textbf{Description} \\
\midrule
$\hat{y}_i, y_i$ & $\mathbb{R}$ & Predicted and ground-truth continuous scalars \\
$e_i$ & $\mathbb{R}$ & Residual error $\hat{y}_i - y_i$ \\
$\delta$ & $\mathbb{R}_{> 0}$ & Huber threshold parameter ($1.0$) \\
$\tau$ & $(0, 1)$ & Target quantile ($0.5$) \\
$\mathcal{L}$ & $\mathbb{R}_{\ge 0}$ & Reduced regression loss \\
\bottomrule
\end{tabular}
\end{table}

\section{Problem Definition and Core Concepts}

\subsection{Intuition and Motivation}
In supervised continuous prediction tasks, the choice of loss function establishes the statistical estimator:
\begin{enumerate}
    \item \textbf{MSE ($L_2$)} corresponds to Gaussian maximum likelihood, estimating the conditional mean $\mathbb{E}[y \mid x]$, but is heavily sensitive to outliers due to quadratic penalization ($e^2$).
    \item \textbf{MAE ($L_1$)} corresponds to Laplace maximum likelihood, estimating the conditional median, and is robust to heavy-tailed noise.
    \item \textbf{Huber Loss} combines the smooth differentiability of MSE near zero ($|e| \le \delta$) with the linear robustness of MAE for large errors ($|e| > \delta$).
    \item \textbf{Quantile Loss} penalizes under-predictions and over-predictions asymmetrically to predict arbitrary conditional quantiles $Q_\tau(y \mid x)$.
\end{enumerate}

\subsection{Formal Mathematical Definition}
\begin{definition}[Loss Formulations and Gradients]
For residual error $e = \hat{y} - y$:
\begin{enumerate}
    \item \textbf{MSE Loss}:
    \begin{equation}
    \ell_{\text{MSE}}(e) = \frac{1}{2} e^2, \quad \frac{\partial \ell}{\partial \hat{y}} = e \label{eq:mse_def}
    \end{equation}
    \item \textbf{MAE Loss}:
    \begin{equation}
    \ell_{\text{MAE}}(e) = |e|, \quad \frac{\partial \ell}{\partial \hat{y}} = \operatorname{sign}(e) \label{eq:mae_def}
    \end{equation}
    \item \textbf{Huber Loss}:
    \begin{equation}
    \ell_{\text{Huber}}(e) = \begin{cases} \frac{1}{2} e^2 & \text{if } |e| \le \delta \\ \delta |e| - \frac{1}{2} \delta^2 & \text{if } |e| > \delta \end{cases}, \quad \frac{\partial \ell}{\partial \hat{y}} = \begin{cases} e & \text{if } |e| \le \delta \\ \delta \operatorname{sign}(e) & \text{if } |e| > \delta \end{cases} \label{eq:huber_def}
    \end{equation}
    \item \textbf{Quantile (Pinball) Loss}:
    \begin{equation}
    \ell_{\text{Quantile}}(e) = \max(\tau e, (\tau - 1) e), \quad \frac{\partial \ell}{\partial \hat{y}} = \begin{cases} \tau & \text{if } e > 0 \\ \tau - 1 & \text{if } e < 0 \end{cases} \label{eq:quantile_def}
    \end{equation}
\end{enumerate}
\end{definition}

\section{Algorithmic Specification}

The computational pipeline implemented in \texttt{NnAlgoRegressionLosses} is shown below.

\begin{algorithm}[H]
\caption{Continuous Regression Loss Evaluation and Gradient Computation}
\begin{algorithmic}[1]
\Require Predictions $\hat{y} \in \mathbb{R}^N$, targets $y \in \mathbb{R}^N$, loss type $\in \{\text{mse}, \text{mae}, \text{huber}, \text{quantile}\}$
\Require Threshold $\delta > 0$, quantile $\tau \in (0, 1)$, reduction $\in \{\text{mean}, \text{sum}, \text{none}\}$
\Ensure Scalar or vector loss $\mathcal{L}$, analytic gradients $g \in \mathbb{R}^N$
\State $N \gets \text{length}(\hat{y})$
\If{$N = 0$ \textbf{or} $\text{length}(y) \ne N$}
    \State \textbf{throw} PreconditionFailureException(``Vector length mismatch or empty'')
\EndIf
\State $\text{losses} \gets \text{new Array of size } N, \quad g \gets \text{new Array of size } N$
\For{$i \gets 1 \textbf{ to } N$}
    \State $e \gets \hat{y}[i] - y[i]$
    \If{$\text{loss\_type} = \text{mse}$}
        \State $\text{losses}[i] \gets 0.5 \cdot e^2, \quad g[i] \gets e$
    \ElsIf{$\text{loss\_type} = \text{mae}$}
        \State $\text{losses}[i] \gets |e|, \quad g[i] \gets \operatorname{sign}(e)$
    \ElsIf{$\text{loss\_type} = \text{huber}$}
        \If{$|e| \le \delta$}
            \State $\text{losses}[i] \gets 0.5 \cdot e^2, \quad g[i] \gets e$
        \Else
            \State $\text{losses}[i] \gets \delta \cdot |e| - 0.5 \cdot \delta^2, \quad g[i] \gets \delta \cdot \operatorname{sign}(e)$
        \EndIf
    \ElsIf{$\text{loss\_type} = \text{quantile}$}
        \State $\text{losses}[i] \gets \max(\tau \cdot e, (\tau - 1.0) \cdot e)$
        \State $g[i] \gets \tau \text{ if } e > 0 \text{ else } (\tau - 1.0 \text{ if } e < 0 \text{ else } 0.0)$
    \EndIf
\EndFor
\State $\text{scale} \gets (1.0 / N \text{ if reduction = mean else } 1.0)$
\State \Return $\{\text{loss}: \sum \text{losses} \cdot \text{scale}, \text{gradients}: g \cdot \text{scale}\}$
\end{algorithmic}
\end{algorithm}

\section{Theoretical Guarantees and Robust M-Estimation}

\begin{theorem}[Huber Loss $\mathcal{C}^1$ Smoothness and Robustness]
The Huber loss function $\ell_{\text{Huber}}(e)$ is continuously differentiable ($\mathcal{C}^1$) everywhere on $\mathbb{R}$, with Lipschitz continuous gradients:
\begin{equation}
|\ell_{\text{Huber}}'(e_1) - \ell_{\text{Huber}}'(e_2)| \le |e_1 - e_2|
\end{equation}
and its influence function $\psi(e) = \ell_{\text{Huber}}'(e)$ is bounded: $\sup_{e \in \mathbb{R}} |\psi(e)| = \delta$, yielding a finite breakdown point resistant to catastrophic outliers.
\end{theorem}

\subsection{Limitations}
\begin{itemize}
    \item \textbf{Hyperparameter $\delta$ Sensitivity}: Choosing $\delta$ too large makes Huber behave like outlier-sensitive MSE; setting $\delta$ too small makes it approximate MAE with slow asymptotic convergence.
\end{itemize}

\section{Complexity Analysis}

\subsection{Time and Space Complexity}
\begin{itemize}
    \item \textbf{Time}: Exactly $\Theta(N)$ arithmetic operations.
    \item \textbf{Space}: $\Theta(N)$ memory for gradient storage.
\end{itemize}

\section{Worked Numerical Example}

Consider $\hat{y} = [3.0, 1.0]^T$, $y = [1.0, 5.0]^T$, errors $e = [2.0, -4.0]^T$, $\delta = 1.0$, $\text{reduction} = \text{mean}$.

\begin{enumerate}
    \item \textbf{Sample 1 ($e_1 = 2.0 > \delta$)}:
    \begin{equation}
    \ell_1 = 1.0(2.0) - 0.5(1.0^2) = 2.0 - 0.5 = 1.50, \quad g_1 = 1.0(+1) = 1.0
    \end{equation}
    \item \textbf{Sample 2 ($e_2 = -4.0$, $|e_2| > \delta$)}:
    \begin{equation}
    \ell_2 = 1.0(4.0) - 0.5(1.0) = 4.0 - 0.5 = 3.50, \quad g_2 = 1.0(-1) = -1.0
    \end{equation}
    \item \textbf{Batch Mean}:
    \begin{equation}
    \mathcal{L} = \frac{1.50 + 3.50}{2} = 2.50, \quad g = [0.50, -0.50]^T
    \end{equation}
\end{enumerate}

\section{Numerical Considerations and Edge Cases}

\begin{itemize}
    \item \textbf{Zero Residual Subgradient in MAE}: At $e=0$, the subgradient is $[-1, 1]$; implementations conventionally set $g = 0.0$.
\end{itemize}

\begin{thebibliography}{9}
\bibitem{huber1964} P.~J.~Huber, ``Robust estimation of a location parameter,'' \emph{The Annals of Mathematical Statistics}, vol.~35, no.~1, pp.~73--101, 1964.
\bibitem{koenker2001} R.~Koenker and K.~F.~Hallock, ``Quantile regression,'' \emph{Journal of Economic Perspectives}, vol.~15, no.~4, pp.~143--156, 2001.
\end{thebibliography}

\end{document}
